% !TeX root = main.tex
\section*{Reviewer 2}

\noindent\textit{Overall assessment.} The reviewer finds that several questions were answered satisfactorily but that the central methodological issues are not resolved: the statistical basis of the variance decomposition, the use of the SABER-FISH data given the averaging requirement, the weaker than stated uncertainty and $\sigma_u$ results, and the tone of the Abstract, Introduction and Discussion. The major points follow.

\bigskip

% ------------------------------------------------------------------ R2.1
\subsection*{R2.1 --- Conditioning variable of the decomposition within one ZT}

\reviewer{In response to my previous comment on the validity of their variance decomposition (R2.2 in the author's response), the authors argue that the conditioning variable $S$ in their Equation 13 (revised main text) is one replicate at one ZT time point and that the $\sigma_{\text{bio}}/\sigma_{\text{tech}}$ split is a ``further decomposition that lives entirely inside the within term''.
While this is OK, the author's decomposition purely within a ZT still must obey the law of total variance, since it is an identity that holds for any conditioning variable.
Therefore, if Eq.~14 is to be more than an analogy, it must itself be a total-variance decomposition within one single ZT timepoint with respect to some conditioning variable, and the authors need to precisely define what that variable is.
Otherwise, Eq.~14 appears quite ad hoc and it is impossible to interpret what the terms $\sigma_{\text{bio}}$ and $\sigma_{\text{tech}}$ are supposed to represent in a statistically rigorous framework.}

\response

% ------------------------------------------------------------------ R2.2
\subsection*{R2.2 --- Consistency with the law of total variance}

\reviewer{In the response R2.2, the authors go on to say that they ``never define $\sigma_{\text{tech}}$ as the within-term of a latent phase-cluster decomposition (an average over heterogeneous sub-populations)\dots\ so there is no averaging over sub-populations with differing variances to reconcile''.
Does this mean the author's procedure is not consistent with the Law of Total Variance, which must have an averaging procedure across sub-populations?
If so, when the authors say that ``At the level of individual time points the additivity is only approximate'', is the approximation coming partly from this departure from the Law of Total Variance?
These points should be more carefully discussed in the Main Text and Methods, since the current version provides little clarity or discussion on this important issue.}

\response

% ------------------------------------------------------------------ R2.3
\subsection*{R2.3 --- Aggregation across the cycle and the SABER-FISH analysis}

\reviewer{The revised Methods state that Eq.~14 ``becomes near-exact ($R^2 \approx 1$) only once variances are aggregated across populations that span the circadian cycle'', and draw from this the requirement that ``the data must cover most of the circadian cycle''.
In the response letter, the authors further comment that ``$\sigma_{\text{bio}}$ is a property of a population sampled across the cycle, not of an individual ZT''.
This is made explicit in Eq.~16 of the revised Main Text.
This then raises the following important questions which I do not currently understand:}

\subsubsection*{R2.3a}

\reviewer{How can this requirement be reconciled with the analysis of the SABER-FISH dataset, where each ZT is analyzed separately to get $\sigma_{\text{bio}}$ and $\sigma_{\text{tech}}$ estimates at each time point?
Given that this is the main dataset with direct experimental validation of the $\sigma_{\text{bio}}$ estimates, the analysis of this data and the limitations should be discussed more carefully.}

\response

\subsubsection*{R2.3b}

\reviewer{The aggregation requirements of the method as well as all the simulations the authors have performed (Fig.~Reb~3A and 3B) seem to require one ground truth $\sigma_{\text{bio}}$ for the entire population of cells across the different ZT time points.
However, this requirement seems directly in conflict with the SABER-FISH data, where $\sigma_{\text{bio}}$ increases with time post Dexamethasone synchronization.
Therefore, the authors need to give a clear justification for applying their method to this dataset.
If this is only a very approximate and empirical application of the method, the authors should accordingly tone down their claims of discovering increasing desynchrony and using this dataset as a strong validation of their method.}

\response

\subsubsection*{R2.3c}

\reviewer{All simulations in the manuscript appear to hold $\sigma_{\text{bio}}$ to one fixed number across ZT time points, which is why the aggregation procedure seems to work via some cancellations.
But if the $\sigma_{\text{bio}}$ per ZT changes, as is the case with the SABER-FISH data, wouldn't this fundamentally create issues in recovering the right $\sigma_{\text{bio}}$ estimates?
Aggregation here would probably recover some sort of average of the individual $\sigma_{\text{bio}}$'s, right?
Therefore in order to claim that scRitmo can extract meaningful increase of desynchrony over time, the authors should at least provide simulations where ground truth $\sigma_{\text{bio}}$ increases across ZTs within a single simulation.
This simulation should then be analyzed with scRitmo and the results benchmarked for accuracy.}

\response

\subsubsection*{R2.3d}

\reviewer{The equations for the aggregation and their corresponding description, both in the main text and response, are quite confusing because key quantities have been vaguely defined.
What exactly do the words ``condition'' and ``batch'' mean here?
What does the term $\sigma_b$ refer to in Eq.~16 of the revised main text?
None of these have been defined.}

\response

% ------------------------------------------------------------------ R2.4
\subsection*{R2.4 --- Aggregation with few ZT samples}

\reviewer{The authors also comment that the aggregation procedure creates a requirement of sufficient ZT samples across 24 hours.
In their simulations, they use 12 and 8 ZT samples, which is greater than the number of samples one would get from typical scRNA-seq datasets.
For example, the Auerbach et al paper which the authors analyze, had only 4 samples spaced 6 hours apart.
To better understand the limits of the aggregation procedure, the authors should demonstrate what happens in such limited datasets --- would accurate estimates of a population $\sigma_{\text{bio}}$ still be possible if there are just 4 samples?}

\response

% ------------------------------------------------------------------ R2.5
\subsection*{R2.5 --- Identical mRNA and protein oscillations in the FLASH-seq data}

\reviewer{In the application of the method to phase inference in asynchronized FLASH-seq data, I appreciate the new results presented and agree that they provide more evidence that biologically meaningful oscillations are being reconstructed.
However, I'm still puzzled by the fact that the mRNA and protein oscillations look identical, and not phase shifted at all.
The protein oscillation peak should be delayed by 2--4 hours as per a vast body of prior literature, which is not reflected in the results here.
How does one interpret this?
Does this suggest that in fact there are significant errors being made in phase assignment to the single cells?
If so, I would recommend discussing this in the Main Text and also toning down the statement ``The inferred phase accurately reconstructed the oscillation of both the endogenous (Rev-Erb$\alpha$) mRNA level and the Rev-Erb$\alpha$-Venus protein fluorescence'', since the meaning of ``accurately'' cannot really be quantitatively established in this setting.}

\response

% ------------------------------------------------------------------ R2.6
\subsection*{R2.6 --- Strength of the $\sigma_u$ claims}

\reviewer{In the main text the authors say ``Across all tissues, cells with higher posterior uncertainty $\sigma_u$ showed larger phase errors than cells with lower $\sigma_u$ (Fig.~S3C, E, F), consistent with the per-cell correlation between $\sigma_u$ and phase error''.
In the Discussion they say ``The posterior uncertainty $\sigma_u$ consistently predicts inference quality across simulations and diverse biological datasets, enabling principled quality control without external time labels.''
However, in the light of the new quantification of correlations between $\sigma_u$ and error (ranging from 0.08 to 0.35 in Fig.~Reb~10A), these statements seem quite disproportionate.
Even the highest of these correlation values at 0.35 is weak at best, and it is also unclear whether any of these would pass statistical tests of significance.
The authors should therefore report these magnitudes along with results from statistical tests in the Main text, and statements such as those above should be tempered accordingly.}

\response
